\documentclass[12pt,letterpaper]{report}
\usepackage{fullpage}
\usepackage{mathtools}
\usepackage[top=2cm, bottom=4.5cm, left=2.5cm, right=2.5cm]{geometry}
\usepackage{amsmath,amsthm,amsfonts,amssymb,amscd}
\usepackage{lastpage}
\usepackage{enumerate}
\usepackage{fancyhdr}
\usepackage{float}
\usepackage{mathrsfs}
\usepackage{xcolor}
\usepackage{graphicx}
\usepackage{listings}
\usepackage{hyperref}
\usepackage{enumitem}
\usepackage{mathtools}
\usepackage{bbm}
\usepackage{amsopn}
\DeclareMathOperator{\Aut}{Aut}


\hypersetup{%
  colorlinks=true,
  linkcolor=blue,
  linkbordercolor={0 0 1}
}

\setlength{\parindent}{0.0in}
\setlength{\parskip}{0.05in}

\newcommand{\CC}{\mathbb{C}}
\newcommand{\NN}{\mathbb{N}}
\newcommand{\QQ}{\mathbb{Q}}
\newcommand{\RR}{\mathbb{R}}
\newcommand{\ZZ}{\mathbb{Z}}
\newcommand{\FF}{\mathbb{F}}
\newcommand{\cO}{\mathcal{O}}


\newtheorem{question}{Question}
%%%%%%%%%%%%%%%%%%%%%%%%%%%%%%%%%%%%%%%%%%%%%%%%%%%%%%%%%% }

%%%%%%%%%%%%%%%%%%%%%%%% CHANGE HERE %%%%%%%%%%%%%%%%%%%% {
\newcommand\course{TECH-GB DATA SCIENCE \& AI FOR BUSINESS}
\newcommand\semester{Spring 2026}
\newcommand{\projecttitle}{Deliverable \#3: Project Update}

%%%%%%%%%%%%%%%%%%%%%%%%%%%%%%%%%%%%%%%%%%%%%%%%%%%%%%%%%% }

%%%%%%%%%%%%%%%%% DO NOT CHANGE HERE %%%%%%%%%%%%%%%%%%%% {
\pagestyle{fancyplain}
\headheight 35pt     
\chead{\textbf{\projecttitle}}
\rhead{\course \\ \semester}
\lfoot{}
\cfoot{\thepage}
\rfoot{}
\headsep 1.5em
%%%%%%%%%%%%%%%%%%%%%%%%%%%%%%%%%%%%%%%%%%%%%%%%%%%%%%%%%% }

\begin{document}
\begin{verbatim}
"""
Authors: 
Cheng(William) Yang - cy2932@nyu.edu
Matthew Cheng - matthewcheng@stern.nyu.edu 
Rebecca KaiWen Yap - ry2678@stern.nyu.edu 

Assignment: Deliverable #3 -- Project Update (Status Report)
Date due: April 11th, 2026
"""
\end{verbatim}

\begin{center}
    \textbf{\large{Deliverable \#3: Project Update (Status Report)}}
\end{center}

\noindent\textbf{Project title:} Profit-Optimized Credit Limit Allocation via Predictive Credit Risk Modeling

\vspace{0.5em}
\noindent\textbf{Team GitHub repository (code, docs, and current progress):} \\
\url{https://github.com/williamyc777/Credit_Limit_Allocation}

\section*{Part 1: Team Members}

\begin{itemize}
    \item \textbf{Cheng (William) Yang} \\
    NetID: cy2932 \\
    Email: \texttt{cy2932@nyu.edu}
    
    \item \textbf{Matthew Cheng} \\
    NetID: mjc9310 \\
    Email: \texttt{matthewcheng@stern.nyu.edu}
    
    \item \textbf{Rebecca KaiWen Yap} \\
    NetID: ry2678 \\
    Email: \texttt{ry2678@stern.nyu.edu}
\end{itemize}

\section*{Part 2: Project Summary (from proposal)}

We study how to support \textbf{credit limit} decisions using LendingClub data: estimate probability of default (PD), treat \textbf{loan amount} as a proxy for credit exposure, and (in later milestones) simulate limit changes and expected profit. The instructor emphasized that PD is \textbf{not fixed} when limits change---our analysis and modeling pipeline are being aligned with this \textbf{prescriptive} perspective (re-scoring under counterfactual \texttt{loan\_amnt}).

\section*{Part 3: Preliminary Results}

\subsection*{3.1 Data pipeline and sample}

\begin{itemize}
    \item \textbf{Source:} Kaggle LendingClub accepted loans (\texttt{accepted\_2007\_to\_2018Q4.csv}).
    \item \textbf{Cleaning:} Retained core features (e.g., \texttt{loan\_amnt}, \texttt{term}, \texttt{int\_rate}, \texttt{grade}, employment length, home ownership, \texttt{annual\_inc}, \texttt{dti}, \texttt{fico\_range\_low}). Restricted to \textbf{terminal} outcomes only: \texttt{Fully Paid} vs.\ \texttt{Charged Off} (excludes \texttt{Current} and other open statuses to avoid label contamination).
    \item \textbf{Label:} Binary \texttt{default} $= 1$ if \texttt{Charged Off}, else $0$.
    \item \textbf{Missing values:} Dropped for the baseline (simple first pass); categoricals one-hot encoded with \texttt{drop\_first=True}.
    \item \textbf{Scale:} After filtering, the cleaned modeling table contains on the order of \textbf{1.27 million} loans; realized default rate approximately \textbf{19--20\%}.
\end{itemize}

\subsection*{3.2 Baseline PD model}

\begin{itemize}
    \item \textbf{Algorithms compared:} Logistic regression, decision tree, random forest, and \textbf{HistGradientBoostingClassifier} (\texttt{scikit-learn}); logistic uses \textbf{standardized} features; best model by hold-out \textbf{ROC-AUC} is saved for scoring and simulation. Test-set \textbf{ROC/PR curve} figures for all models and a \textbf{reliability (calibration)} plot for the AUC-best model are written by \texttt{src/evaluation\_plots.py} at the end of training.
    \item \textbf{Split:} Train/test (e.g., 80/20), stratified on \texttt{default}, fixed random seed for reproducibility.
    \item \textbf{Outputs:} \texttt{clean\_data.csv}, \texttt{pd\_predictions.csv} (per-model PD columns + \texttt{PD} aligned to best model), \texttt{best\_model.pkl}, \texttt{all\_models.pkl}, \texttt{scaler.pkl}, \texttt{feature\_cols.pkl}, \texttt{model\_comparison.csv} (\textbf{AUC} and \textbf{PR-AUC} for imbalanced defaults).
    \item \textbf{Ranking performance:} Report both ROC-AUC and PR-AUC from \texttt{train\_model.py} (insert latest numbers from \texttt{model\_comparison.csv}).
    \item \textbf{Decision logic (instructor feedback):} We do \textbf{not} treat a \textbf{0.5} cutoff on PD as a business rule---it is only sklearn's default for hard labels from \texttt{.predict()}. We document \textbf{explicit decision analysis} via \texttt{src/decision\_logic.py} (grid of candidate thresholds under a simple retrospective \textbf{profit proxy} on the test split), and we rely primarily on \textbf{continuous PD} for ranking and for upcoming \textbf{limit / profit simulation}.
\end{itemize}

\subsection*{3.3 Exploratory analysis and visualization}

\begin{itemize}
    \item \textbf{Notebook:} \texttt{creditallocation\_eda.ipynb} documents business-oriented EDA on the same \texttt{clean\_data} and \texttt{pd\_predictions} files.
    \item \textbf{Key exhibits:} FICO vs.\ default; \textbf{loan amount} binned by quantiles (deciles) with default rates by bin; interest rate and DTI vs.\ default; default rates by grade, term, and home ownership (derived from one-hot columns); scatter of predicted PD vs.\ \texttt{loan\_amnt}; ROC curve for predicted PD vs.\ realized default.
    \item \textbf{Narrative:} A dedicated section, \textit{Implications for Credit Limit Optimization}, links EDA to the prescriptive goal: \texttt{loan\_amnt} as exposure proxy, why PD should change when limits change, and connection to simulation below.
\end{itemize}

\subsection*{3.4 Prescriptive simulation and threshold analysis}

\begin{itemize}
    \item \texttt{src/decision\_logic.py} --- scan approval thresholds on \texttt{P(default)} under a simple \textbf{retrospective profit proxy} (instructor: do not center analysis on 0.5).
    \item \texttt{src/simulation\_profit.py} --- grid of \textbf{counterfactual} \texttt{loan\_amnt} multipliers; \textbf{re-score PD} with the best model for each alternative; choose the multiplier that maximizes a stylized \textbf{one-period expected profit} per loan; export portfolio and sample-level summaries.
    \item \texttt{src/causal\_pd\_analysis.py} --- small ``what-if'' demo (change amount, observe PD) using the saved best model.
\end{itemize}

\subsection*{3.5 Collaboration, repository, and progress to date}

The team maintains a \textbf{public GitHub repository} for reproducibility and shared development:

\begin{center}
\url{https://github.com/williamyc777/Credit_Limit_Allocation}
\end{center}

\noindent\textbf{What is already in the repository (reflects work completed as of this update):}

\begin{itemize}
    \item \textbf{Documentation:} \texttt{README.md} describes the business problem, baseline pipeline, optional output bundle (\texttt{model\_outputs\_bundle.zip}), and how EDA ties to credit-limit optimization.
    \item \textbf{Data cleaning and labels:} \texttt{src/preprocess.py} --- feature selection, terminal-status filter, \texttt{default} label, one-hot encoding, writes \texttt{output/clean\_data.csv} (file gitignored locally; teammates regenerate or use bundle).
    \item \textbf{PD models:} \texttt{src/train\_model.py} --- four models (including gradient boosting), AUC and PR-AUC, \texttt{best\_model.pkl} and \texttt{pd\_predictions.csv}; \texttt{src/evaluation\_plots.py} --- ROC/PR and calibration figures under \texttt{output/}.
    \item \textbf{Supporting scripts:} \texttt{src/eda\_pd\_distribution.py}; \texttt{src/causal\_pd\_analysis.py}; \texttt{src/decision\_logic.py}; \texttt{src/simulation\_profit.py}; \texttt{run\_pipeline.sh} (auto-selects Anaconda Python when available); \texttt{package\_output.sh}; \texttt{docs/COMPLETE\_PROJECT\_GUIDE.md}; \texttt{PROJECT\_STATUS.md} (checklist).
    \item \textbf{Exploratory analysis:} \texttt{creditallocation\_eda.ipynb} --- business-oriented visuals (FICO, DTI, rate, binned \texttt{loan\_amnt} vs.\ default rate, PD vs.\ amount, ROC) and narrative on prescriptive use of \texttt{loan\_amnt}.
    \item \textbf{Dependencies:} \texttt{requirements.txt} for Python packages.
\end{itemize}

\noindent\textbf{Repository hygiene:} Large Kaggle CSVs and most \texttt{output/*} artifacts are \textbf{gitignored} to stay within GitHub limits; a \textbf{zipped output bundle} may be committed when small enough so collaborators can run the notebook without retraining.

\noindent\textbf{Further iteration (optional):} XGBoost/LightGBM, calibration plots, fairness discussion, production deployment---out of scope for the core class deliverable but listed in \texttt{PROJECT\_STATUS.md}.

\section*{Part 4: Analysis of Preliminary Results}

\subsection*{4.1 Substantive patterns}

\begin{itemize}
    \item \textbf{Credit quality:} Lower FICO (range low) aligns with higher default frequency, consistent with underwriting intuition.
    \item \textbf{Pricing and capacity:} Higher interest rates and higher DTI mass among defaulters support treating pricing and repayment capacity as co-predictors with exposure.
    \item \textbf{Exposure (\texttt{loan\_amnt}):} Binned default rates show how risk co-moves with loan size in historical data---a descriptive prior for why \textbf{perturbing loan amount in the model} should move predicted PD, matching instructor feedback on prescriptive analytics.
    \item \textbf{Model alignment:} PD vs.\ \texttt{loan\_amnt} and ROC analysis indicate the models provide a sensible \textbf{ranking} signal; \textbf{operating points} (approve/decline or limit changes) should be chosen by \textbf{stated business logic}---not an arbitrary 0.5 probability cutoff.
\end{itemize}

\subsection*{4.2 Implications for the next modeling stage}

\begin{itemize}
    \item The \textbf{profit objective} will combine predicted PD with revenue and loss assumptions; the baseline gives a first PD surface to \textbf{re-evaluate under counterfactual limits} using the saved model and scaler.
    \item The \textbf{prescriptive path} (re-score under counterfactual \texttt{loan\_amnt}, expected profit on a grid) is implemented in code. We may add \texttt{class\_weight}, richer profit accounting, or optional \textbf{XGBoost/LightGBM} (beyond sklearn HGB) as refinements for the final report.
\end{itemize}

\section*{Part 5: Issues and Challenges}

\begin{itemize}
    \item \textbf{Class imbalance:} High accuracy on a default 0.5 label rule is misleading; we focus on \textbf{PD as a continuous input}, \textbf{threshold / profit analysis} where appropriate, and \textbf{simulation-based} policy comparison rather than a single hard cutoff.
    \item \textbf{Causal interpretation:} Historical correlation of loan size and default reflects \textbf{selection} (who gets larger loans) and macro factors; we treat the logistic model as a \textbf{scenario engine} for ``what-if'' limits, not as a randomized experiment.
    \item \textbf{Data and tooling:} Large file sizes complicate GitHub workflows; teammates must regenerate \texttt{output/} locally or use shared bundles. Notebook paths must be run from the repository root so \texttt{output/} resolves correctly.
    \item \textbf{Refinement:} Tune stylized \texttt{margin}/\texttt{lgd} and revenue assumptions in \texttt{simulation\_profit.py} to match the report's business story; optional extra learners (XGBoost/LightGBM) or additional calibration if time permits; incorporate professor's Chapter~11-style discussion of \textbf{single-model counterfactuals} in the final write-up.
    \item \textbf{Optional environment note:} Some laptops showed architecture mismatches for scientific Python stacks; the team standardizes on a consistent conda/venv where needed.
\end{itemize}

\section*{Part 6: Next steps (brief)}

\begin{enumerate}[label=\arabic*.]
    \item Integrate \textbf{simulation and threshold} results (CSV summaries from \texttt{simulation\_profit.py} and \texttt{decision\_logic.py}) into the final report and presentation figures.
    \item Align \textbf{profit parameters} and narrative with the business case; use \texttt{pr\_curves\_all\_models.png} and other \texttt{output/} figures in the final deck (ROC/PR/calibration are already produced by the pipeline).
    \item Final polish: executive summary, limitations, and team-specific contributions.
\end{enumerate}

\end{document}
